% Author: GPT-6 Astra Pro, /root/replacement, 2026-11-22.
% Numerical effort not exposed. Continued from HODGE_DIVISION checkpoint15.
\section{Follow-up 1: the semiregular replacement and its exact source maps}
\label{sec:f1-replacement}

This section belongs to DEG\_L Q-T7-REPLACEMENT and the existing
HODGE\_DIVISION T5-EIGHTFOLD row. The actual split-sixfold
construction HD-N01--HD-N09 remains a proved background result.

\subsection{HD-F1-R01: the scope of S-30--S-36 and the actual comparison}

The objects constrained in DEG\_L S-30--S-36 live on an abelian
fourfold $X$ with Rosati-invariant real multiplication by
$F=\mathbf Q(\sqrt{\Delta_0})$, $\Delta_0>1$ squarefree, of
type $(2,2)$. Their specified full character is
\[
 e^{-c_1(E)/N}\operatorname{ch}(E)=Nv_q,
 \qquad v_q=1-qC+q^2\nu,
 \quad C=\frac{\theta_1^2+\theta_2^2}{2},
 \quad\nu=\frac{\theta^4}{24},\quad N>0,\quad q>0.
 \tag{F1R1.1}
\]
Here $N$ is the virtual rank on $X$.
The integral-normalizer exclusion in S-30 additionally assumes
ordinary simplicity, gluability and full semiregularity.
S-35's conclusion $\Delta_0\mid N$ additionally assumes a simple
gluable fully semiregular equivariant lift. S-31 is a vector-bundle
assertion. The local maximal-order, principal-polarization and
determinant-space assumptions in S-32, S-33 and S-36 are retained.

For an abelian $m$-fold $X$, the actual secant-square functor is
\[
 \Phi=(\mathrm{id}_X\boxtimes\mathrm{FM}_{\mathcal P^{-1}[m]})
          \mu^*,\qquad \mu(x,y)=(x+y,y).
\]
Grothendieck--Riemann--Roch, with all Todd factors one, gives
\[
 \Phi_H(v_1\boxtimes v_2)=(-1)^m(p_{x,\alpha})_*
 \left(\mu^*(v_1\boxtimes v_2)e^{-c_1(\mathcal P_{y,\alpha})}\right).
 \tag{F1R1.2}
\]
Restriction to a fibre over $(x,\alpha)$ consequently gives
\[
 \operatorname{rk}\Phi(E_1\boxtimes E_2)
       =(-1)^m\int_X\operatorname{ch}(E_1)\operatorname{ch}(E_2).
 \tag{F1R1.3}
\]
Translation and a degree-zero line bundle leave the cohomological
integrand fixed. This proves the formula directly. For two characters
$N_iv_{q_i}$ on a fourfold, its value is
$DN_1N_2(q_1+q_2)^2$, where $D=\int_X\nu$.
Thus the ranks in S-30--S-36 have a specified source-to-target map.

In Markman's real-quadratic construction, the quartic CM field
$K=F(\sqrt{-q})$ acts on $A=X\times\widehat X$. With
$K_0=\mathbf Q(\sqrt{-q})$, the corresponding Weil spaces are
\[
 W_K=\bigwedge_K^4H^1(A,\mathbf Q)\subset H^4(A,\mathbf Q),
 \qquad
 W_{K_0}=\bigwedge_{K_0}^8H^1(A,\mathbf Q)\subset H^8(A,\mathbf Q).
\]
For each embedding $\tau:K_0\to\mathbf C$, let $\sigma_1,\sigma_2$
be its two extensions to $K$. There is the exact exterior map
\[
 \bigwedge^4H^1_{\sigma_1}(A)\otimes
 \bigwedge^4H^1_{\sigma_2}(A)
 \xrightarrow{\ \wedge\ }
 \bigwedge^8\left(H^1_{\sigma_1}(A)\oplus
                         H^1_{\sigma_2}(A)\right).
 \tag{F1R1.4}
\]
It is an isomorphism of lines. Together with its conjugate it gives,
after the rational $K_0$-Weil projector, a surjection
\[
                    \operatorname{Sym}^2_{\mathbf Q}W_K
                              \longrightarrow W_{K_0}.
 \tag{F1R1.5}
\]
The projector is the endomorphism polynomial selecting the two
endpoint eigenspaces in $H^8$, as in DEG\_L B6. The two line maps
prove surjectivity after extension to $\mathbf C$ and hence over
$\mathbf Q$.

The already constructed sixfold control has its own exact morphisms.
HD-N03 uses threefold quotient ideals; HD-N05 gives
\[
 S=\begin{pmatrix}1&1\\2i_0&-2i_0\end{pmatrix}^{\oplus3},
 \quad JS=S\operatorname{diag}(i_0,-i_0),\quad\deg S=4096.
\]
The actual idempotents in HD-N06 give
$\operatorname{ch}(F)=S^*\operatorname{ch}(\mathcal E)/16$,
$\operatorname{rk}(F)=4$, and $\sigma_{S^*\mathcal E}=16\sigma_F$
on the induced invariant obstruction classes. The specified final
line-bundle tensor produces a weakly equivariantly semiregular
complex on the balanced CM sixfold with group order4096 and
\[
 \operatorname{ch}(\mathcal E_d)=32t_d+4-8\theta_0^2
                 +\frac83\theta_0^4-\frac{16}{45}\theta_0^6,
 \qquad B_0=0.
 \tag{F1R1.6}
\]
This is an existing S-14 construction. It also refutes the unqualified
intake sentence excluding integral normalizers for every S-14
object: its refined replacement is S-30 with the ordinary simple
gluable balanced-RM fourfold assumptions written above.

\subsection{HD-F1-R02: a complete fourfold point-orbit benchmark}

Let $X=E_i^2$ with its product principal polarization $H$. Choose
$T=\{0,h\}\subset X[2]$ of order two and let $r:X\to Y=X/T$.
Define the actual perfect complex on $A=X\times\widehat X$ by
\[
 F_0=r^*I_{0,Y}=I_T,\quad L=\mathcal O_X(H),\quad F=F_0\otimes L,
 \qquad
 E=(\mathrm{id}\boxtimes\mathrm{FM}_{\mathcal P^{-1}[2]})
                        \mu^*(F\boxtimes F).
 \tag{F1R2.1}
\]
It has a coherent finite point action of order four, no negative
self-Ext, and injective degree-two trace on its invariant self-Ext.
Its whole twisted character is
\[
               e^{-P/2}\operatorname{ch}(E)
                       =t-2\theta+\frac{\theta^3}{12},
 \tag{F1R2.2}
\]
where $\theta$ is the product principal polarization on $A$,
$P=c_1(\mathcal P)$, and $t$ is primitive integral and nonzero in
the middle Weil plane for the balanced $\mathbf Q(i)$ action below.
In particular this is a rank-zero S-14 object, with
\[
                  \theta t=0,\qquad\int_A t^2=8,
                  \qquad\int_A\theta^4/24=1.
 \tag{F1R2.3}
\]

\paragraph{The trace and action proof.}
On the abelian surface $Y$, $I=I_0$ is simple: every endomorphism
extends to a scalar on its double dual $\mathcal O_Y$. Serre
duality makes $\operatorname{Ext}^2(I,I)\to H^2(\mathcal O_Y)$
the dual of the identity map on $\operatorname{Hom}(I,I)$, hence
an isomorphism. Since $\operatorname{ch}(I)=1-[0]$, Riemann--Roch
gives $\chi(I,I)=-2$ and
\[
                  (\operatorname{ext}^0,\operatorname{ext}^1,
                    \operatorname{ext}^2)(I,I)=(1,4,1).
\]
Degree-one contraction with $1-[0]$ has rank four: the two
$H^1(\mathcal O_Y)$ directions wedge the scalar to independent
classes, while the two tangent directions contract the point class
to independent elements of $H^2(\Omega_Y^1)$. Evaluation followed
by trace is this contraction. Since the middle self-Ext space has
dimension four, both evaluation and trace are injective. Degree-zero
trace sends the identity to $1-[0]$ and is injective.

For $I\boxtimes I$, the three second-self-Ext summands have
dimensions $1,16,1$. The trace is, up to the tensor sign,
\[
 \sigma_I^2\otimes\operatorname{ch}(I),\quad
 \sigma_I^1\otimes\sigma_I^1,\quad
 \operatorname{ch}(I)\otimes\sigma_I^2
 \tag{F1R2.4}
\]
into the distinct Hochschild bidegrees $(-2,0),(-1,-1),(0,-2)$.
Every factor is injective. Finite-etale descent identifies
\[
 \operatorname{Ext}^j(F_0,F_0)^T
                      =\operatorname{Ext}^j(I,I),
\]
and pullback on Hochschild homology is an isomorphism intertwining
trace. The line-bundle tensor, $\mu^*$ and Fourier equivalence in
(F1R2.1) conjugate the entire coherent action and its trace maps.
The original product is a sheaf; equivalence preserves its negative
self-Ext vanishing.

Identify $\widehat X$ with $X$ by $L$. Before applying $\mu^*$,
$(h_1,h_2)\in T^2$ acts by translation $(h_1,h_2)$ and tensor
coordinate $(-h_1,-h_2)$. After $\mu^*$ these become
$(h_1-h_2,h_2)$ and $(-h_1,-h_1-h_2)$.
The inverse Fourier kernel sends translation $a$ to tensor $P_a$
and tensor $P_b$ to translation $-b$. The final subgroup is
\[
 G=\{((h_1-h_2,h_1+h_2),(-h_1,h_2)):h_1,h_2\in T\}
                              \subset A\times\widehat A.
 \tag{F1R2.5}
\]
Its tensor coordinates recover both $h_i$, so it has order four.
Its coherence is inherited from descent. The invariant trace
injection just proved gives weak $G$-semiregularity through the
criterion used in DEG\_L S-14.

\paragraph{The whole character calculation.}
Use integral coordinates $a_j,b_j,u_j,v_j$, $j=1,2$, and put
\[
 \theta=\sum_j(a_jb_j+u_jv_j),\quad P=\sum_j(a_jv_j-b_ju_j),
 \quad\Omega_{s,j}=a_jb_j+s(a_jv_j-b_ju_j)-u_jv_j,
\]
\[
 t=\tfrac12(\Omega_{i,1}\Omega_{i,2}
                        +\Omega_{-i,1}\Omega_{-i,2}).
\]
The input character equals
\[
 \operatorname{ch}(F)=1+H-H^2/2
       =\alpha e^{iH}+\beta e^{-iH},\qquad
 \alpha=(1-i)/2,\quad\beta=(1+i)/2.
\]
For one elliptic pair, the equal-sign Fourier integral, after
the factor $e^{-P/2}$, is $2s\exp(s\theta_j/2)$; the opposite-sign
integral is $\Omega_{s,j}$. Multiplying the two pair factors with
the input coefficients gives $-4\sin(\theta/2)+t$. Its expansion
is exactly (F1R2.2). The rational exterior-algebra certificate
\texttt{check\_surface\_seed.py} independently expands all twelve
generators, integrates the four fibre variables, and verifies this
identity, primitivity, and all three equalities (F1R2.3).

On each elliptic pair the integral action
\[
                   J=\begin{pmatrix}0&1\\-1&0\end{pmatrix}
\]
satisfies $J^2=-1$, is skew-adjoint for the product polarization,
and has the two opposite holomorphic eigenvalues. Its cotangent
eigenvectors are $u_j-sa_j,v_j-sb_j$, and
$\Omega_{s,j}=-(u_j-sa_j)(v_j-sb_j)$.
The two products defining $t$ therefore lie on the two Weil
determinant lines. Their displayed integral monomials verify
$t\ne0$. This proves every claimed S-14 datum at the fourfold.
The construction is a small-dimensional control in the secant-square
design, with no originality claim.

\paragraph{The exact dimension-dependent step.}
The surface shortcut used here is the duality map
\[
 \operatorname{Ext}^2_X(I,I)\longrightarrow H^2(\mathcal O_X)
 \quad\text{dual to}\quad
 H^0(\mathcal O_X)\longrightarrow\operatorname{Hom}_X(I,I).
 \tag{F1R2.6}
\]
On an abelian surface simplicity makes the latter an
isomorphism. Markman's published ICM2026 paper, Section2,
paragraph defining the semiregularity map, states precisely
this surface argument. For a threefold Serre duality instead
identifies $\operatorname{Ext}^2(I,I)$ with
$\operatorname{Ext}^1(I,I)^*$, so (F1R2.6) alone supplies
no corresponding conclusion. The successful replacement is
already proved in HD-N02: the plane-quartic canonical map
$\operatorname{Sym}^2H^0(K_C)\to H^0(K_C^2)$ is an
isomorphism, giving invariant Ext dimensions $(1,6,6,1)$;
contraction with $1-h^2+2h^3/3$ has full degree-two rank6.
The external-product and Fourier maps of HD-N03 then give
the actual sixfold construction in HD-N06. Thus there is
no step proved to fail for every $n\ge3$.

There is an exact failure for the simple gluable ample
imaginary-secant seed on a fourfold, already HD-N11--HD-N12.
Its full character is $e^B(\alpha e^{sH}+\bar\alpha e^{-sH})$,
$s^2=-d<0$, and its contraction ranks are $(8,12)$.
For an equivariantly simple gluable lift the full trace
injection would give $e_2=12$, $e_1\ge8$ by the invariant
CY4 transpose map, and hence
\[
 \chi_G=2-2e_1+e_2\le-2,
 \qquad
 \chi_G=\frac{32d^2\alpha\bar\alpha}{|G|}
                     \int_XH^4/24>0,
 \tag{F1R2.7}
\]
a contradiction. This is the precise excluded seed step;
the balanced-RM and general coupled-object continuations
retain their separately stated scopes.

\subsection{HD-F1-R03: balanced-RM products and the quadratic eightfold character}

Consider an external product of two actual fourfold objects.
One normalized full character is of balanced-RM form $Nv_q$ as in
(F1R1.1). The other has either that form or the genuine ample
imaginary two-exponential form in HD-N11. Allow rational algebraic
$B$-field factors on both characters. No object in the orbit of
this product under Fourier--Mukai equivalences, finite isogeny
pullbacks or pushforwards, line-bundle tensors, translations,
shifts and derived duals has full character
\[
 e^{B_A}\bigl(P(\theta_A)+\lambda t_A\bigr),\quad\lambda\ne0,
 \tag{F1R3.1}
\]
on a quadratic-Weil eightfold of signature $(4,4)$, for a compatible
ample $\theta_A$ and a nonzero middle Weil class $t_A$.
This exclusion needs no simplicity or semiregularity assumption.

Indeed write $c_v^j:HH^j\to HH_{-j}$ for contraction with the
whole character. DEG\_L S-22 gives ranks $(8,18)$ for the balanced
character; HD-N11 gives $(8,12)$ for the imaginary character.
For the first pair, the two real-multiplication complex eigenspaces
give two factors of contraction ranks $(4,1)$. Their sum and
Hochschild--Kunneth maps give $8$ and $1+16+1=18$, without
assuming a rational product decomposition of $X$.

On the external product the actual second contraction is
\[
 (c_{v_1}^2\otimes v_2)\ \oplus\
 (c_{v_1}^1\otimes c_{v_2}^1)\ \oplus\
 (v_1\otimes c_{v_2}^2).
 \tag{F1R3.2}
\]
The three target Hochschild bidegrees are distinct. The exact
second ranks are therefore $18+64+18=100$ and $18+64+12=94$.

For every actual equivalence the intertwiner is
\[
 \Phi_{HH}(c_v^j(a))
      =c_{\Phi_{HH}(v)}^j(\Phi^{HH}(a)),
 \tag{F1R3.3}
\]
with both transport maps invertible. Isogeny pullback has the
same identity through its differential; pushforward is
$(\deg f)(f^*)^{-1}$. A line-bundle tensor gives the invertible
shear $(v,w)\mapsto(v,w+\iota_vB)$; translations act identically,
shifts multiply by a sign, and dualization conjugates the degree
signs. Thus the displayed ranks are preserved throughout the
stated functor orbit.

The exact whole-character calculation HD-N14 gives instead
\[
 \operatorname{rank}c^2=56+16\rho\in\{56,72,88,104\},
 \qquad\rho=\operatorname{rank}\langle P,P',P''\rangle\bmod z^7
 \tag{F1R3.4}
\]
for every character (F1R3.1). The two ranks94 and100 are absent,
proving the assertion. In particular the external square of
HD-F1-R02, after undoing its specified maps, has rank100 as the
product of four surface point ideals. This recovers the earlier
HD-N14/HD-R05 obstruction through that explicit construction.

This closes the stated balanced-RM product supplier, with no
claim of nonexistence for arbitrary coupled eightfold complexes.
The relation (F1R1.5) between quartic and quadratic Weil spaces
remains the explicit multiplication map; no semiregular object
with the quadratic whole character is constructed from it here.

\subsection{HD-F1-R04: the surviving nonsimple RM lift has many endomorphisms}

Keep (F1R1.1), a specified coherent action of an effective finite
$G\subset X\times\widehat X$, and an actual equivariant lift of
$E$. Suppose its invariant negative self-Ext groups vanish and
the full ordinary trace is injective on $\operatorname{Ext}^2(E,E)^G$.
Simplicity is not assumed. Define
\[
 e_j=\dim\operatorname{Ext}^j(E,E)^G,\quad
 k=\frac{4DN^2q^2}{|G|},\quad
 \Delta_{\mathrm{out}}=\prod_{\substack{p\mid\Delta_0\\p\nmid N}}p.
\]
S-14 requires vanishing of all ordinary negative self-Ext
groups. The present necessary calculation uses only the
invariant vanishing, so applies in particular to those S-14
lifts; these numerical assumptions alone are not an object
construction for S-14.
Then every such lift satisfies
\[
 e_2=18,\quad e_1\ge8,\quad
 k=2e_0-2e_1+18\in2\mathbf Z_{>0},\quad e_0\ge k/2-1,
 \tag{F1R4.1}
\]
and
\[
 \begin{array}{c|c|c}
 &\text{divisibility}&\text{endomorphism bound}\\\hline
 N\text{ odd}&4\Delta_{\mathrm{out}}\mid k
             &e_0\ge2\Delta_{\mathrm{out}}-1\\
 N\text{ even}&2\Delta_{\mathrm{out}}\mid k
             &e_0\ge\Delta_{\mathrm{out}}-1.
 \end{array}
 \tag{F1R4.2}
\]
In particular a rank-one lift has $e_0\ge2\Delta_0-1$;
invariant endomorphism dimensions one and two are excluded.
For rank one, $\Delta_0=2$ and $e_0=3$ force the profile
$(e_0,e_1,e_2,e_3,e_4)=(3,8,18,8,3)$.
This specifies a necessary profile, without an existence claim.

\paragraph{Proof.}
The coefficient-sector proof in S-21 supplies the identity-sector
Hochschild isomorphism, Serre functor $[4]$ and exact Euler factor
$\chi_G=\chi_X/|G|$. It uses actual translation graphs and nontrivial
flat line bundles, and no simplicity hypothesis. The invariant
degree-two evaluation map is the Serre transpose of trace. Trace
injectivity makes evaluation surjective, so their composite has
rank $e_2$. S-22 computes that rank as18 and the degree-one rank
as8. Invariant Serre duality gives $e_4=e_0$, $e_3=e_1$, and the
assumed negative-Ext vanishing removes all other degrees.
Riemann--Roch gives $\chi_X(E,E)=4DN^2q^2>0$.
This proves (F1R4.1).

Let $H$ be the product of the primary subgroups of $G$ at primes
not dividing $N$, and put $G_N=G/H$. The nonsimple descent
assertion of S-35 applies to this actual lift and gives
\[
 \frac{N^4Dq^2}{|H|}
       =\frac{N^2k|G_N|}{4}\in\Delta_0\mathbf Z_{>0}.
 \tag{F1R4.3}
\]
Every prime in $|G_N|$ divides $N$. Each odd prime in
$\Delta_{\mathrm{out}}$ consequently divides $k$.
For odd $N$, the factors $N$ and $|G_N|$ are odd, so
integrality forces $4\mid k$, and $2\mid\Delta_0$ forces
$8\mid k$. This gives $4\Delta_{\mathrm{out}}\mid k$.
For even $N$, $\Delta_{\mathrm{out}}$ is odd; combining its
divisibility with the evenness of $k$ gives
$2\Delta_{\mathrm{out}}\mid k$. The positive lower bounds
on $k$ and (F1R4.1) prove (F1R4.2). The endpoint profile follows
from $8\le k\le8$ and invariant Serre duality.

\paragraph{The stronger reduced-rank form.}
Let $T=\ker(G\to X)$, let $R$ be the radical of its scalar
commutator, and put
\[
 r=|R|,\quad d=\sqrt{|T|/r},\quad m=N/(rd),\quad
 \Delta_m=\prod_{\substack{p\mid\Delta_0\\p\nmid m}}p.
\]
S-38 gives positive integers $d,m$. Without any split Brauer
assumption, the stronger bounds are
\[
 \begin{array}{c|c|c}
 &\text{divisibility}&\text{endomorphism bound}\\\hline
 m\text{ odd}&4\Delta_m\mid k&e_0\ge2\Delta_m-1\\
 m\text{ even}&2\Delta_m\mid k&e_0\ge\Delta_m-1.
 \end{array}
 \tag{F1R4.4}
\]
Here $e_0$ is still measured on the original equivariant lift.

Indeed retain S-41's actual center isogeny $p:B\to X$ of degree
$r$, residual Brauer class $\beta$ of period $a\mid m$, cover
$u=[a]:B'\to B$, and the prime-to-$m$ subgroup $H_0$ of
$H=G/T$. Its translations lift by $h\mapsto cb_h$, where
$ac\equiv1\pmod{\exp H_0}$. The transported coherent action
and ordinary Morita transform give an actual $H_0$-equivariant
perfect object $F'$ of rank $m$ with
\[
 e^{-c_1(F')/m}\operatorname{ch}(F')=m u^*p^*v_q,
 \qquad D'=a^8rD.
\]
These constructions do not use the specialization $k\in\{2,4\}$.
Writing $H_m=H/H_0$, the S-35 descent relation on $F'$ is
\[
 \frac{m^4D'q^2}{|H_0|}
   =\frac{a^8m^2 k|H_m|}{4}
       \in\Delta_0\mathbf Z_{>0}.
 \tag{F1R4.5}
\]
All prime divisors of $a$ and $|H_m|$ divide $m$.
The same valuation argument as for (F1R4.3) proves (F1R4.4).
No semiregularity of $F'$ or identification of its self-Ext
groups with those of the original lift is asserted. In particular
reduced rank one requires $e_0\ge2\Delta_0-1$, also for $N>1$.

\paragraph{Literature and remaining scope.}
Mukai, Nagoya Math. J.81 (1981), Theorem2.2 and formula(3.1),
supplies the precise Fourier equivalence and translation/tensor
identities. Caldararu--Willerton, New York J. Math.16 (2010),
Sections4--7, supplies the functorial trace framework.
Markman's arXiv:2502.03415v2 and arXiv:2509.23079v1 are design
and comparison sources, with no unrefereed theorem used as a
premise. His ICM 2026 paper, DOI10.1137/25M1803796,
Sections2 and11, confirms the small-dimensional scope of that
design. The current search included arXiv, Hexagon and journals.
The close quartic-CM rank exclusions in Bhattacharjee--Bhattacharya,
\emph{Secant Objects for Weil Classes on Abelian Varieties},
\url{https://www.preprints.org/manuscript/202611.0672},
Theorems41--42, are recorded as comparison, not as proof inputs.
No originality claim is made for the integration consequences here.
The one-fibre existential S-14 task in a dimension $2n$ with
$n\ge3$ is already proved by HD-N06 at $n=3$, whose actual
maps and whole character are displayed in HD-F1-R01. The
unqualified transfer of S-30--S-36 to every target object is
corrected by the exact source hypotheses and rank map of
HD-F1-R01. The balanced-RM seed and an actual coupled
eightfold object are separate open continuation questions.
